\chapter{Conclusion}

\section{Summary of Key Findings}

This study estimated the causal effect of ITN use on malaria infection among children under five in Kenya using KMIS 2015 and 2020 data (6,771 children, 528 clusters). Key findings:

\begin{enumerate}
	\item \textbf{Confounding is substantial}: Unadjusted models showed harmful association (OR $\approx$ 1.1-1.2); fully adjusted models revealed protective effect (OR = 0.592 for KMIS 2015).
	
	\item \textbf{Clustering matters}: ICC of 41-47\% justifies random effects models \citep{killip2004intracluster}.
	
	\item \textbf{GLMM-based causal methods show strong protective effects}: PSM, IPTW, and DR using GLMM for propensity scores produced ORs of 0.446-0.554 (45-55\% reduction). As \citet{neuhaus1991comparison} note, these conditional effects are larger than marginal effects when ICC is high.
	
	\item \textbf{DML provides conservative estimates}: DML produced ORs of 0.961-0.981 (2-4\% reduction), representing population-average effects.
	
	\item \textbf{Heterogeneity exists}: Strongest effects in poorer wealth quintile, urban areas, older children (36-59 months), and high-rainfall regions.
\end{enumerate}

\section{Answer to the Research Question}

Based on triangulation across six causal methods, ITN use has a protective causal effect on malaria infection among children under five in Kenya. Cluster-specific effects show 45-55\% reduction in malaria odds; population-average effects show 2-4\% reduction. At the population level, a 2.5\% reduction translates to approximately 125,000 malaria cases averted annually.

\section{Main Contribution}

\subsection{Methodological Contribution}

This study demonstrates the value of combining traditional causal methods with modern machine learning for triangulation. The use of GLMM for propensity scores to account for hierarchical structure (ICC = 41-47\%) is an important advancement. Following \citet{ayele2024assessment}, we extended multilevel modeling to propensity score estimation. The balancing property of the propensity score \citep{rosenbaum1983central} ensures covariate balance. The DML framework \citep{chernozhukov2018double} provided flexible, non-parametric estimation of population-average effects.

\subsection{Policy Contribution}

This study provides rigorous causal evidence that ITNs reduce malaria infection. Heterogeneity analysis identifies subgroups where ITNs are most effective, informing targeted distribution strategies. The decline in ITN use from 58.0\% (2015) to 50.5\% (2020) highlights the need for behavior change communication.

\section{Future Research Directions}

Future research should address: (1) long-term effects as nets age, (2) combination interventions (ITNs + IRS + vaccines), (3) mechanisms explaining why ITNs are less effective for the poorest, (4) generalizability to other countries, (5) insecticide resistance monitoring, (6) DML variance estimation following \citet{wu2025why}, (7) positivity assessment methods \citep{westreich2010invited}, and (8) longitudinal designs.

\section{Final Remarks}

This thesis demonstrates that ITN use reduces malaria infection among children under five in Kenya. The triangulation of evidence across six causal methods provides robust support. The distinction between cluster-specific effects (45-55\% reduction) and population-average effects (2-4\% reduction) highlights the importance of understanding which estimand is relevant for different policy questions.

From a policy perspective, the findings support continued investment in ITN distribution programs, with attention to targeting and complementary interventions for the most vulnerable. The decline in ITN use from 2015 to 2020 is concerning; behavior change communication is needed.

Methodologically, this study illustrates the value of modern causal machine learning methods for estimating treatment effects from observational data.

In summary, ITNs work. They reduce malaria infection in children. But they work better for some children than others. Targeting ITN distribution to high-risk areas and populations, combined with complementary interventions for the most vulnerable, can maximize the impact of malaria control programs in Kenya and beyond.